\chapter{Introduzione}
\label{cap:introduzione}

\begin{definizione}
  Un vettore dello spazio si scrive nella base canonica come
  $\vec{v} = a\,\ihat + b\,\jhat + c\,\khat$.
\end{definizione}

\begin{teorema}[Norma euclidea]
  \label{teo:norma}
  Per ogni $\vec{v} = a\,\ihat + b\,\jhat + c\,\khat$ si ha
  \[
    \lVert \vec{v} \rVert = \sqrt{a^2 + b^2 + c^2}.
  \]
\end{teorema}

\begin{proof}
  Applicare due volte il teorema di Pitagora.
\end{proof}

\begin{esempio}
  Se $\vec{v} = 3\,\ihat + 4\,\jhat$, allora $\lVert \vec{v} \rVert = 5$
  (teorema~\ref{teo:norma}).
\end{esempio}

\begin{osservazione}
  Aggiungi altri capitoli nella cartella \texttt{capitoli/} e includili con \verb|\input|.
\end{osservazione}
